\chapter{Data Formats}
\label{ch:data-formats}

\section{Colour}
\label{sec:colour}

A colour is a 16-bit word with 5 bits for each of red, green and blue. Bit 15 is a reserved flag.

\begin{figure}[H]
	\centering
	\begin{bytefield}[bitwidth=1.5em]{16}
		\bitheader{0-15} \\
		\bitbox{1}{F} &
		\bitbox{5}{R} &
		\bitbox{5}{G} &
		\bitbox{5}{B}
	\end{bytefield}
	\caption{Colour Format}
	\label{fig:colour}
\end{figure}

Each 5-bit channel drives the converter of its colour output (Chapter~\ref{ch:package}), which has 32 evenly spaced
levels from black to full scale. Channel value 31 is full scale, so a pixel of 31 in every channel is full-scale white.

\section{The Palette}
\label{sec:palette}

The palette holds 256 colours at offsets \texttt{0x6000}--\texttt{0x60FF} (Chapter~\ref{ch:memory-map}). Every pixel of
every background is a palette lookup. Palette entry 0 is special in one way: it is the backdrop, the colour shown where
no layer has a visible pixel.

\section{Tile Data}
\label{sec:tile-data}

A tile is a small square image of 4-bit pixels, each of which is a number from 0 to 15. The value 0 is transparent: the
pixel shows whatever is behind it. The values 1--15 select a colour through the palette
(Section~\ref{sec:palette-select}).

\subsection{Mode 1 Tile Layout}
\label{sec:m1-tile-layout}

Four pixels fill one 16-bit word, with the leftmost pixel in the most significant nibble.

\begin{figure}[H]
	\centering
	\begin{bytefield}[bitwidth=1.5em]{16}
		\bitheader{0-15} \\
		\bitbox{4}{P1} &
		\bitbox{4}{P2} &
		\bitbox{4}{P3} &
		\bitbox{4}{P4}
	\end{bytefield}
	\caption{Four Pixels of Tile Data}
	\label{fig:tile-word}
\end{figure}

An 8 by 8 tile is stored as eight rows of two words each, top row first, so it occupies 16 words. The word for pixels
$4h$ to $4h + 3$ of row $r$ of tile $t$ is at the tile data base address plus $16t + 2r + h$, where $h$ is 0 for the
left half of the row and 1 for the right half. A 16 by 16 tile is described in Section~\ref{sec:flat-16}.

\subsection{Mode 2 Tile Layout}

Mode 2 stores a 16 by 16 texel tile in a different arrangement, described in Section~\ref{sec:m2-tile-format}. It takes
64 words, the same as a 16 by 16 tile would in the row-by-row layout, but the words are not rows.

\subsection{16 by 16 Orthogonal Tiles}
\label{sec:flat-16}

An orthogonal background whose tile size bit (Section~\ref{sec:reg-tilesize}) is 1 uses 16 by 16 tiles. A 16 by 16 tile
is stored as four 8 by 8 quadrants in the order top left, top right, bottom left, bottom right, each laid out as in
Section~\ref{sec:m1-tile-layout}, so a tile is 64 words. For the pixel in column $c$ and row $r$ of tile $t$ the word
is at

\begin{center}
	\texttt{tile data base + 64 $\times$ tile number + 16 $\times$ quadrant + 2 $\times$ (row mod 8) + half}
\end{center}

where the quadrant is $2\lfloor r/8 \rfloor + \lfloor c/8 \rfloor$ and the half is $\lfloor (c \bmod 8)/4 \rfloor$. A
2-bit 16 by 16 tile is four quadrants of 8 words, 32 words, with the word at \texttt{base + 32 $\times$ tile + 8
	$\times$ quadrant + (row mod 8)}. A tile of 64 words is the same size as an isometric tile, so one tile number
addresses either. With 16 by 16 tiles a 32 by 32 tilemap covers 512 by 512 pixels. Each background chooses its own tile
size, and backgrounds of different tile sizes need no further rule: each is fetched independently.

\subsection{Flipping Orthogonal Tiles}

When the \texttt{H} bit of an orthogonal tilemap entry is set, the tile is mirrored left to right: pixel column $c$ is
read from column $w - 1 - c$ of the tile, where $w$ is 8 or 16. When \texttt{V} is set, row $r$ is read from row $w - 1
	- r$. The device reads the mirrored word and reverses the order of its pixels.

\subsection{2-Bit Orthogonal Tiles}
\label{sec:flat-two-bit}

A background whose \texttt{T2} bit (Section~\ref{sec:reg-palette}) is set uses 2-bit tiles. An 8 by 8 orthogonal tile
is then 8 words, one for each row, holding the row's eight pixels with the leftmost in bits 15--14:

\begin{center}
	\texttt{tile data base + 8 $\times$ tile number + row}
\end{center}

Value 0 is transparent, so a pixel has three colours. A word covers eight pixels, so the device reads it for every
second group of four. 2-bit tiles apply to every orthogonal background and to the status layer. Isometric 2-bit tiles
have a layout of their own (Section~\ref{sec:m2-two-bit}).

\section{Tilemap Entries}
\label{sec:tilemap-entry}

A tilemap is a table of one-word entries, 32 entries to a row and 32 rows, stored row by row. An entry names a tile and
how it is drawn.

\begin{figure}[H]
	\centering
	\begin{bytefield}[bitwidth=1.5em]{16}
		\bitheader{0-15} \\
		\bitbox{1}{V} &
		\bitbox{1}{H} &
		\bitbox{1}{P} &
		\bitbox{3}{PAL} &
		\bitbox{10}{TILE}
	\end{bytefield}
	\caption{Tilemap Entry}
	\label{fig:tilemap-entry}
\end{figure}

\begin{table}[H]
	\centering
	\begin{tabular}{llp{7.5cm}}
		\toprule
		\textbf{Bits} & \textbf{Field} & \textbf{Meaning}                                                                                                                                     \\
		\midrule
		15            & \texttt{V}     & Flip the tile vertically on an orthogonal layer (Section~\ref{sec:flat-16}); darkens the tile in an isometric layer (Section~\ref{sec:m2-shading})   \\
		14            & \texttt{H}     & Flip the tile horizontally on an orthogonal layer; lightens the tile in an isometric layer                                                           \\
		13            & \texttt{P}     & Priority: 1 puts the tile in front of everything, sprite layer 1 included; 0 leaves its layer in the order S1, B1, S2, B2, S3, B3 (and B4 in Mode 1) \\
		12--10        & \texttt{PAL}   & Palette selection, 0 to 7 (Section~\ref{sec:palette-select})                                                                                         \\
		9--0          & \texttt{TILE}  & Tile number, 0 to 1023, counted from the background's tile data base address                                                                         \\
		\bottomrule
	\end{tabular}
	\caption{Tilemap Entry Fields}
	\label{tab:tilemap-entry-fields}
\end{table}

Tile numbers can reach 1023, but VRAM is not that large for every tile size: 64-word tiles stop at 512 in a VRAM with
nothing else in it (Chapter~\ref{ch:memory-map}). A tile number whose data lies past the end of the address space wraps
at 16 bits.

\section{Tilemap Sizes}
\label{sec:tilemap-size}

A tilemap is made of screens of 32 by 32 entries, stored row by row, 1,024 words each. The width bit \texttt{BnX} and
the height bit \texttt{BnY} of the Tile Size register (Section~\ref{sec:reg-tilesize}) double the map across or down:

\begin{center}
	\begin{tabular}{ccl}
		\toprule
		\texttt{BnX} & \texttt{BnY} & Map                                                 \\
		\midrule
		0            & 0            & One screen, 32 by 32 entries                        \\
		1            & 0            & Two screens side by side, 64 by 32: A B             \\
		0            & 1            & Two screens one above the other, 32 by 64: A over B \\
		1            & 1            & Four screens, 64 by 64: A B over C D                \\
		\bottomrule
	\end{tabular}
\end{center}

The screens are in memory in the order A, B, C, D, starting at the tilemap base. For the entry in column $c$ and row
$r$ of a map $W$ entries wide and $H$ high (both after wrapping, $c \bmod W$ and $r \bmod H$), the address is

\begin{center}
	\texttt{tilemap base + 1024 $\times$ screen + 32 $\times$ (r mod 32) + (c mod 32)}
\end{center}

where the screen is $\lfloor r/32 \rfloor \cdot (W/32) + \lfloor c/32 \rfloor$. A map wraps at its size. The tilemap
sizes apply to orthogonal backgrounds, where an entry is a tile, and to isometric backgrounds, where an entry is a cell
(Section~\ref{sec:m2-tilemap}).

\section{Palette Selection}
\label{sec:palette-select}

A background pixel with a non-zero value $v$ in a tile whose entry has palette selection $s$ reads palette entry
\[
	\text{entry} = ((s \ll k) + \texttt{OFFSET} + v) \bmod 256,
\]
where \texttt{OFFSET} and $k$ come from the background's palette configuration register
(Section~\ref{sec:reg-palette}): $k$ is 2, 3, 4 or 5 for \texttt{SIZE} values 0, 1, 2 and 3, that is, for palettes of
4, 8, 16 or 32 colours. The sum wraps at 8 bits. A pixel whose value is 0 is transparent and reads no palette entry.

With 16-colour palettes, \texttt{OFFSET} 0 and a selection of 3, the value 5 reads palette entry 53. The eight
selections therefore pick eight palettes of 16 colours, 128 palette entries in all, starting at the offset.

\section{Compositing Rule}
\label{sec:compositing}

Layers are combined pixel by pixel. For each screen pixel the device takes the first layer, in front-to-back order,
whose pixel value is not 0, and shows the palette entry that pixel selects. If every layer is transparent, the device
shows palette entry 0. There is no blending: a pixel is either transparent or fully opaque.

The front-to-back order of the layers is S1, B1, S2, B2, S3, B3, B4 in Mode 1, the same without B4 in Mode 3, and S1,
B1, S2, B2, S3 with a status layer in Mode 2; see Sections~\ref{sec:m1-priority} and~\ref{sec:m2-layers}. A tile with
its priority bit set is in front of everything.
